% M. Carlsen - F.Caruana, Vugar Gashimov Memorial 2014: your comments. Format: games/README.md.
@ title | M. Carlsen - F.Caruana, Vugar Gashimov Memorial 2014

@ 6...O-O | board
\mythoughts{Magnus took the pawn on move 6; he could equally have castled with 6.O-O.
I am unsure whether 7.b4 is sound here---White has an extra tempo in a Catalan-style
setup. In the main Catalan Black often posts the kingside bishop on e7; here White's stands on g2 rather than e2, so it is the more active piece.

I prefer 7.O-O.}

@ 7...a5
\mythoughts{8.b4 axb4 9.cxb4 \N e4 is bad: White drops a rook.

In Ruy Lopez or Benoni lines where Black has pawns on a6, b5 and c4, a knight on d3 often
cramps White.

White's structure is analogous: the extra pawn sits on c5, and a knight on d6 would be tempting if Black pushes ...e5 and seizes the center.

Black's last move also weakens b5. I would like 8.\N a3.}

\myreflection{8.\N a3 is playable. Then I play 9.\B e3 transpires.}

@ 8...Nc6
\mythoughts{I still want 9.\N a3, developing the queenside.}

@ 9...a4
\mythoughts{I am not sure what Black wants from ...a4. Perhaps \N a5-c4, eyeing the loose b2-pawn,
or a trade against the bishop on e3 that spoils White's structure.

10.\N d4 e5 11.\N xc6 bxc6 leaves Black with a firm center and an open b-file.

10.\N b5 is premature: the d6 outpost is not yet secure.

Black stands well in the center; crude lunges with h4 or f4 are unlikely to work.

Karpov has observed that sometimes doing nothing is fine: you let the opponent act,
and their plans may create weaknesses you can exploit.

I would happily wait, but I must still choose a move.
10.h3 worries me because it softens g3. If Black builds ...e4, striking the center with f3 will be awkward while the g3-pawn remains tender.

Black may try 10...\N g4 to gain a tempo, then 11...f5.
10.\Q c1 prevents that idea and looks strong.}

\myreflection{10.\N d4 e5 11.\N db5 \B e6 12.\N d6 e4 13.\N ab5 \N e8 14.\N xe8 \R xe8 15.b4 White is better. I don't have to take the knight on c6 as I have calculated.

10.\N b5 \N g4 is indeed premature. The knight on b5 doesn't have a target.

10.h3 e5 11.\N b5 White still has an advantage.My worry about weakening g3 was unsound.}

@ 10...e5 | board
\mythoughts{Black intends 11...d4 12.cxd4 exd4 13.\B f4 \R e8 14.\R e1 \N e4, regaining the pawn with the more active army.

I dislike 11.\N b5 \Q a5 12.\N d6 \N d7: Black recovers the pawn and keeps the livelier pieces.

11.\R d1 addresses the point. I keep ideas such as \B g5, \N xe5, and \R xd5 in reserve.}

\myreflection{Black actually doesn't threat with 11...d4?! 12.\R d1 since the pawn is pinned.

Also 11.\N b5 \Q a5?! 12.\N d6 \N d7? 13.\R d1 White has an advantage.

I realize now how I miss the ``obvious'' move \R d1 in the other variations, which
I exactly planned as a first move. I don't know why this happens. Certainly it
is something to be aware of.}

@ 11...Qe7
\mythoughts{12.\B h6 \B xh6 13.\Q xh6 \Q xc5 is poor: Black seizes firmer control of the center.

12.c4 d4 13.\N xd4 exd4 14.\B xd4 \N xd4 15.\R xd4 also occurs to me,
but I distrust it after 15...\Q xe2, 15...\B e6, and similar annoyances.
I need not decide which reply is most painful---that is the spirit of alpha-beta pruning.

12.\N b5, eyeing d6, looks tempting. Even from a passive shell I would get some counterplay.}

\myreflection{12.\B h6 is indeed bad and my calculation is correct.

12.c4?! d4 13.\N xd4? is indeed poor. After 13.\B f4 \Q xc5, Black has some slight advantage.}

@ 12...Be6
\mythoughts{Black plans 13...\R a5 14.\N d6 \N c5. After ...\N xc5, the light-squared bishop has left c8; White's knight on d6 may be trapped, while the bishop on e6 still covers d5.

Black may also try 13...\N d7 14.\N d6 f5.

13.\N g5 h6 is awkward for White: trading the bishop opens the f-file,
and retreating the knight squanders two tempi.

13.\N d6 looks strong.}

\myreflection{The assessment is incorrect. 13.\N g5 h6? 14.\N xe6 fxe6 15.c4 White
has some advantage. Opening the f-file is not promissing for Black since he
has no target there.

13.\N d6 \N g4 14.\N g5 is indeed good.}

@ 13...Bg4
\mythoughts{I still want 14.\N d6.}

@ 14...h6
\mythoughts{15.\N f3 looks natural.}

@ 15...Kh7
\mythoughts{Black aims for 16...e4 17.\N d2 \B xe2, when his game looks more comfortable;
hence 16.h3.}

@ 16...Be6
\mythoughts{17.b4 stirs counterplay; 17.\K h2 is equally reasonable.
Order hardly matters---Black cannot prevent both.

I would choose 17.\K h2.}

\myreflection{17.\K h2?! \N e8 is too passive. Black starts a counterstrike with 18...f5 while
White was just waiting.}

@ 19...Ne4
\mythoughts{I did not see ...Ne4 coming. Capturing on e4 feels limp:
20.\N xe4 dxe4 21.\N e1 f5 22.\N e1 \R d8 23.\R xd8 \Q xd8. Both queens
command open files, and Black should drum up enough counterplay.

20.b4 \N d8 21.\N xe4 dxe4 22.\N e1 f5 leaves me stuck on the queenside while Black's center still offers resources.

20.c4 \N xd6 21.cxd6 \Q xd6 22.cxd5 \B xd5 23.\B xh6 \B xh6 24.e4
regains the piece and sets the queenside pawns rolling. Black no longer enjoys the same central clamp as in the earlier lines, so White should keep an edge.

I play 20.c4.}

\myreflection{20.c4 \N xd6 21.cxd6 \Q xd6? 22.cxd5 \B xd5? 23.\R xd5 \Q xd5 24.\N g5+ hxg5 25.\B xd5 White is winning.
20.c4 d4! 21.\B xd4 \N xg3 22.fxg3 \N xd4 23.\N xd4 exd4 24.\Q a7 \B xh3 25.\B xh3 \Q xe2 Black is piece
down but has enough counterplay because White's king is exposed.

The move b5 doesn't exist!}

@ 20...f5
\mythoughts{Black now has concrete threats:
21...f4 22.gxf4 exf4 23.\B d4 \B xd4 24.cxd4 \Q h4, after which he should be winning before long.

White must watch f2 and h3. Trading the knight and tucking the king to h2 covers both ideas.

I take 21.\N 2xe4.}

@ 21...dxe4 | board
\mythoughts{The b3-pawn is loose, and 22...f4 must be answered.

22.\Q b1 looks forced.}

@ 22...f4
\mythoughts{I dare not allow 23...f3: Black can break through with ...h5-h4.
Even if I seal the kingside with g4, a sacrifice on that square remains hard to rule out.

Thus 23.gxf4 is obligatory.}

\myreflection{23.gxf4 exf4 24.\B d2 f3 25.exf3 exf3 26.\B f1 \B h3 White lost the advantage. Ironically, it happens exactly as I have feared.

23.\B d2 f3? 24.\B f1 fxe2 25.\B xe2 \B xh3 26.\N xe4 Black has no more attack.

I was panicking too much. When facing a scary attack, it is
helpful to work out what the exact threat was. Only then can
I have a calm mind.}

@ 23...e3
\mythoughts{Black has a scary attack. I dislike 24.\B e1 fxg3 because
the pawn on e3 clamps my position.

I suspect Black intends 24...e4 then 25...\B e5.

24.fxe3 fxg3 25.\B e1 wins the pawn on g3. My bishops shelter the king,
and I may yet hold.

I decide on 24.fxe3.}

\myreflection{24.fxe3 fxg3 25.\B e1 \Q h4 26.\N e4 White still has advantage.

My fear was unrealistic. 24.\B e1 fxg3 25.fxg3 h5 26.b4 Black has no further attack.
Using computer it is clear that Black has no attack. I need to work out without a computer when
defending.}

@ 24...Bf5
\mythoughts{25.\N xf5 gxf5 gives Black an open g-file, which feels too dangerous.
After 25.\B e4 \B xh3 26.\B xg6 \K h8, my king is too exposed.
25.\Q c1 \Q e6 26.\N xf5 gxf5 27.\R d6 \Q e7 28.\Q d1 may offer some counterplay.

I decide on 25.\Q c1.}

\myreflection{After 25.\N xf5 gxf5 26.g4 \Q xc5 27.gxf5 \K h8 28.f6 \R xf6 29.\Q d3 e4 30.\B xe4, Black has no open g-file as I feared.

Even after 25.\B e4 \B xh3 26.\B xg6 \K h8 27.f3 White also has advantage.}

@ 25...h5
\mythoughts{Black intends ...h4. I decide on 26.fxe3 to shore up against that thrust.}

@ 26...fxg3
\mythoughts{I must take on g3; otherwise Black creates a protected passed pawn after ...h4.
I might survive the middlegame, but not the resulting endgame.}

@ 27...Qg5
\mythoughts{After retreating the bishop from g3, I will lose the h3-pawn.
28.e4 \Q xg3 29.exf5 gxf5: Black crashes through on the g-file, and I have no satisfactory defense.
28.\N xf5 is my only try.}

\myreflection{The comment regarding e4 was exaggerated because White is able to attack even though
Black has an open g-file.

28.\N xf5 gxf5 29.\B h2 is indeed good.

Amazingly, all the variations allows Black to open the g-file. What was I scared of?}

@ 28...Qxg3 | board
\myreflection{I miss the text move, an intermediate move.}

@ 30...gxf5
\mythoughts{Black intends 31...\N e7; the knight joins the attack via g6 and f4.
I want 31.\B xc6 bxc6 32.\Q e3, angling for \Q g3 and \R d7 for counterplay.}

\myreflection{31.\B xc6 e4! suddenly Black is winning. It makes no sense to trade my important defender for a
knight that doesn't do too much.}

@ 31...fxe4
\mythoughts{32.\B xe4+ followed by 33.\R f3, and White should be winning.}

@ 33...Rf4 | board
\myreflection{my previous 33.\R f3 \R xf3 34.\B xf3 \Q xh3 35.\B g2 \Q g4 loses a pawn. The
position would be a draw.

In defending I am eager to trade pieces, missing sometimes simple tactics.}

\mythoughts{I feel I have an advantage. I prefer not to trade queens, since
endgames with opposite-colored bishops favor the defender. 34.\N f5 \Q f6 35.\R d7 and then
36.\Q g3 should be decisive.}

\myreflection{34.\N f5? \Q g5+ White loses material.}

@ 34...Qe7
\mythoughts{My worst piece is the rook on d3. I want to improve it with 35.\R d2, heading for a2.}

\myreflection{35.\R d2 \B h6 36.\R f2 \Q f6 37.\R xf4 \B xf4 38.\Q d3 White should be indeed winning.}

@ 35...Qh4
\mythoughts{The rook is still my worst piece, so I offer an exchange with 36.\R f3.
After 36...\R xf3 37.\Q xf3, my structure improves and my king is safer.}

\myreflection{The idea is in principle correct. I have indeed big material advantage and
will be able to exchange the rooks. According to computer, I am still winning. I am however not sure if I have a clear plan
to win the game because opposite-colored bishops favors the attacker and
clearly Black is the attacker.}

@ 36...e4
\mythoughts{37.\N xe4 \B e5 38.\R f3 \R xf3 39.\Q xf3 \Q e1: Black should have a perpetual.
37.\R e3 \N d4 38.cxd4 \B xd4 39.\K h1 \B xe3 40.\Q xe3 \R f3: I see no mating net while Black threatens a perpetual.
37.\R d5 \N e5: I do not see a good follow-up.
I want 37.\R d1 intending 38.\R f1. Trading rooks on f1 is
cleaner than on f3, since it does not open checks against my king.}

\myreflection{37.\N xe4 \B e5 38.\R f3? leads indeed to perpetual check. Instead White should
play 38.\R d5. I have considered the move as candidate and here I forgot it completely,
as if the move only exist as a first move, not as a follow-up move.

37.\R e3 \N d4 38.cxd4 \B xd4 39.\K h1 \B xe3 40.\Q xe3 \R f3 41.\Q h6+ \K g8 42.\Q g6+ \K f8 43.\B xf3 \Q xh3+ 44.\K g1 White is winning. True, White doesn't have a mating threat. I was unable to see
some intermediate moves to win material for White.

37.\R d5 \N e5 38.\N xe4 White is winning. The same topic comes again and again. It seems that I always
fail to combine moves. A move in another line can certainly be used. Why did I miss it so often?}

@ 37...Ne5 | board
\mythoughts{38.\R g3 \R xe4 drops material.
I must reroute the rook. Black's attack has lost its immediate sting.
38.\R d1 followed by 39.\R f1 is a natural exchange of rooks.
I cannot simply retreat the knight from e4 because Black has 40...\Q g3.
The counterattack along the dark squares is awkward.

After 38.\R d5, I do not see a clear plan.

I want 38.\R e3 intending 39.\N d6. Once the b7-pawn
falls, I can advance the queenside pawns for the full point.}

\myreflection{The idea regarding 38.\R e3 is OK.}

@ 38...Kg8
\mythoughts{Now I see the idea: 39.b5 then 40.c6. The h5-pawn should fall
if Black captures on c6.}

@ 39...Rf5
\mythoughts{Black intends 40...\N f3. 40.\K h1 is very natural.}

@ 40...bxc6
\mythoughts{I now see the tactical point if 40...\N f3 41.\B xf3 \R xd5 \N f6, winning the rook.
41.bxc6 is natural.}

@ 41...Qe7
\mythoughts{42.\Q a2 threatens a discovered check, but after 42...\K h8 I achieve little.

I dislike letting the rook slip to the back rank, so I try 42.\N g3. The rook cannot
retreat because the h5-pawn hangs. After 42...\R g5, however, 43.\Q a2 fails:
the knight hangs.

Refining further, I consider 42.\N d6. After 42...\R g5 43.\Q a2 \K h8 44.\Q a8+ \K h7,
45.\B e4 is unavailable because the bishop is pinned.

45.\K h1 still looks strong; I do not see a defense for Black.

I decide on 42.\N d6.}

\myreflection{I like the process. Interestingly, computer says the text move is not precise.
My first candidate 42.\Q a2 \K h8 43.\K h1 is the best. I don't fully understand why.}

@ 42...Rg5
\mythoughts{My earlier calculation was wrong: 43.\Q a2 \N f3 loses for White. If I can force c7, however,
the picture improves.
So I decide on 43.\N b5.}

@ 43...Qe6
\mythoughts{My rook is loose because the bishop is pinned. I play 44.\K f1 to
break the pin.}

\myreflection{There is no need to play a waiting move, if a concrete move is available.}

@ 44...Kh7
\mythoughts{I still would love to play 45.\K h1 because I did't want to calculate in detail
what would happen when Black takes my h3 pawn.
I was not sure about 45.\Q e4+ \R g6, until I found 46.\R d6}

\myreflection{My idea 45.\Q e4+ \R g6 46.\R d6 exchanges the material and wins. Magnus played a more direct one.}

@ 46...Qa6
\mythoughts{Black threatens 47...\Q a1+ 48.\K h2 \N f3+ 49.\Q xf3 \B e5 with mate to follow.
47.\R a8 parries the attack, after which 48.c8=Q decides.}

\myreflection{I like my idea, quite simple. If I am winning, I don't want to check if
my opponent has a perpetual check. Exchanging the queens is the most effective way.}

@ 49.Ke1 | board
\myreflection{Beyond the concrete blunders, this game shows that I still need to improve my defensive technique.
In several moments I reacted to Black's attacking ideas too emotionally, and that led me to overestimate
the danger. I should first identify the exact threat, then compare candidate defenses calmly.

I also noticed a recurring calculation weakness: I can find a good candidate move as the first move in one
line, but I fail to recognize the same move as a follow-up move in a related variation. This breaks the
consistency of my analysis and makes me miss strong continuations. I need to train this specifically:
when I evaluate candidate moves, I should actively check whether they can reappear one move later in
neighboring lines.}

